\documentclass[a4paper,10pt]{article}
%\usepackage[latin1]{inputenc} % Paquetes de idioma (otro encoding)
\usepackage[utf8]{inputenc} % Paquetes de idioma
\usepackage[spanish]{babel} % Paquetes de idioma
\usepackage{graphicx} % Paquete para ingresar gráficos
\usepackage{grffile}
\usepackage{hyperref}
\usepackage{fancybox}
\usepackage{amsmath}
\usepackage{amsfonts}
\usepackage{listings}
% Paquetes de macros de Circuitos
%\usepackage{pstricks}
\usepackage{tikz}

% Encabezado y Pié de página
\usepackage{fancyhdr} % Paquete para encabezados y pie de página
\pagestyle{fancy} % Sin esta línea no se imprimiría el encabezado en todas las páginas

\fancyhf{} %  Borra el encabezado anterior (Por defecto escribe el títutlo de la sección en la que se encuentra la hoja
\setlength{\headheight}{22.55pt}
\fancyhead[L]{
	{\textsf{Facultad de Ingenier\'ia $-$ Universidad de Buenos Aires \\ 66.09 Laboratorio de Microcomputadoras}}
}
%\addtocounter{page}{5}
\fancyhead[R]{\thepage}

\renewcommand{\footrulewidth}{0.4pt} % Ajusta el tamaño de las líneas separadoras en el pié de página
\renewcommand{\headrulewidth}{0.4pt} % Ajusta el tamaño de las líneas separadoras en el encabezado

\fancyfoot[L]{
	{\textsf{Anteproyecto} \\
	{\textsf{Integrantes: Torres Feyuk, Levi Hadid}}
	}
}
		

% Carátula del Trabajo
\title{ \author{} % Lo pongo para que el warning no moleste :p
\setlength{\unitlength}{1cm} %  Especifica la unidad de trabajo
\thispagestyle{empty}

\begin{picture}(18,0)
\put(0,0){\includegraphics[width=1.5cm, height=3cm]{Imagenes/Logo1.png}}

\put(10.5,0){\includegraphics[width=3cm, height=3cm]{Imagenes/Logo2.png}}

\end{picture}
\\[1.5cm]
\begin{center}
	\textbf{{\Huge Facultad de Ingenier\'ia \\ Universidad de Buenos Aires}}\\[2cm]
	{66.09 Laboratorio de Microcomputadoras}\\[0.5cm]
	{Anteproyecto}\\[2.5cm]
\end{center}

\begin{flushleft}
	\textbf{Integrantes:} \\[1cm]

	\begin{tabular}{|c|c|c|}
		\hline
		\textbf{\normalsize Padr\'on} & \textbf{\normalsize Nombre} & \textbf{\normalsize Email} \\
		\hline
		\normalsize 89579 & \normalsize Torres Feyuk, Nicol\'as R. Ezequiel & \normalsize ezequiel.torresfeyuk@gmail.com \\
		\hline
		\normalsize 90406 & \normalsize Levi Hadid, Lucas Alberto & \normalsize lucaslh9@hotmail.com \\
%		\hline
%		\normalsize ????? & \normalsize Madariaga, Eduardo & \normalsize madariagaedu@gmail.com \\
		\hline
	\end{tabular}
\end{flushleft}
\date{} % Hace que no se imprima la fecha en la cual se compilo el .tex
 }

\begin{document}
	\maketitle % Hace que el título anterior sea el principal del documento
	\newpage

	\tableofcontents % Esta línea genera un indice a partir de las secciones y subsecciones creadas en el documento
	\newpage

	\section{Objetivos}
		El objetivo del presente trabajo práctico consiste en diseñar e implementar un sistema que permita controlar los motores de un autito RC mediante 
		comunicación inalámbrica. 

	\section{Prestaciones Técnicas}
		\begin{itemize}
			\item El autito debe contar con 4 velocidades de funcionamiento, algo similar a una caja de cambios.
			\item La comunicación inalámbrica debe abarcar un rango de al menos 10 metros. 
			\item Se debe tener un control absoluto del auto a través del control remoto
		\end{itemize} 

	\section{Descripción del Proyecto}
		El proyecto consiste en realizar el sistema de control de un auto de juguete el cual utiliza motores de continua para desplazarse, de forma tal de manejar 
		al mismo de forma remota a través de un control. Dado que lo que se desea es diseñar y construir el sistema y no la parte mecánica, se utiliza en el presente
		proyecto  el chasis y los motores de un autito de juguete. El autito a utilizar posee dos motores de continua en la parte trasera del mismo, por lo cual
		el mismo posee tracción diferencial. 	
	
		
	\section{Diagramas en Bloque}
	\section{Circuitos Esquemáticos}
	\section{Listado de Componentes}
	\section{Software}
	\section{Resultados}
	\section{Conclusiones}

\end{document}
